\section{Career Evidence: Further Results}\label{app:careers}

\subsection{Industries}

Industry codes of all PSID waves and of the CPS Outgoing Rotation Groups are mapped to the 222 industries of the
ind1990ddx scheme and grouped into 33 industry groups.\footnote{1970 Census codes are bridged to 1980 codes with the
Census Bureau's 1970--1980 table, weighted by the male experienced civilian labor force; 2000 codes use Dorn's weighted
crosswalk; 2002 and 2012 codes use the first three digits. Code: \texttt{IGE/32.IndustryCareers.py}
(\texttt{output/industry\_careers.md}).} For each industry $j$, among men aged 25--64 in the 1992--2002 CPS, I estimate
its wage level at 40, $P_j$, and the slope of its age--wage profile, $B_j$, relative to the average industry
(Table~\ref{tab:industrygroups}). Sons of richer parents sort into better-paid and steeper industries by 25--34
(Table~\ref{tab:industrysorting}). Writing the son's log wage as his industry's level, plus its career slope times his
age, plus a within-industry residual, the career slope accounts for 0.084$^{***}$ of the pooled rise of wage
persistence and 0.051$^{***}$ of the change in the racial gap (Table~\ref{tab:industrydecomp}); with each industry's
full nonlinear profile, 0.118$^{***}$ and 0.074$^{***}$, and with the industry fixed at 25--34, 0.132$^{***}$ and
0.086$^{***}$ (Table~\ref{tab:industrynonlinear}). Same-employer wage growth rises with parental income for White sons
within detailed industries (0.015$^{***}$) and most steeply in high-paying service industries (0.040$^{**}$), and not for
Black sons (Table~\ref{tab:industrygrowth}). Following the father's industry is described in
Table~\ref{tab:fatherindustry}.

\begin{table}[!htbp]
\centering
\caption{Industry wage levels and career slopes (CPS MORG men 25--64, 1992--2002)}
\label{tab:industrygroups}
\resizebox{\ifdim\width>\textwidth\textwidth\else\width\fi}{!}{%
\resizebox*{!}{\ifdim\totalheight>0.85\textheight 0.85\textheight\else\totalheight\fi}{%
\begin{tabular}{lccccc}
\toprule
Industry group & \begin{tabular}[b]{@{}c@{}}Level $P$\\at 40\end{tabular} & \begin{tabular}[b]{@{}c@{}}Career slope $B$\\per 10 years\end{tabular} & \begin{tabular}[b]{@{}c@{}}Raw growth\\25--34 to\\45--55\end{tabular} & \begin{tabular}[b]{@{}c@{}}Return to\\experience\\5 to 25 years\end{tabular} & \begin{tabular}[b]{@{}c@{}}Employment\\share\end{tabular} \\
 & (1) & (2) & (3) & (4) & (5) \\
\midrule
Education & 0.082 & 0.077 & 0.432 & 0.437 & 0.055 \\
Legal services & 0.538 & 0.069 & 0.340 & 0.358 & 0.006 \\
Professional services (engineering, accounting, research, consulting) & 0.394 & 0.039 & 0.314 & 0.380 & 0.024 \\
Manufacturing: chemicals, petroleum, rubber, plastics & 0.115 & 0.038 & 0.314 & 0.340 & 0.029 \\
Other health services & 0.093 & 0.031 & 0.317 & 0.264 & 0.015 \\
Hospitals & 0.071 & 0.029 & 0.314 & 0.485 & 0.022 \\
Communications & 0.254 & 0.022 & 0.297 & 0.412 & 0.018 \\
Manufacturing: motor vehicles, aircraft, other transport eq. & 0.156 & 0.022 & 0.279 & 0.331 & 0.033 \\
Finance: banking, credit, securities & 0.395 & 0.021 & 0.322 & 0.387 & 0.022 \\
Insurance & 0.320 & 0.020 & 0.315 & 0.393 & 0.014 \\
Manufacturing: instruments, misc. manufacturing & 0.097 & 0.020 & 0.301 & 0.337 & 0.014 \\
Manufacturing: textiles, apparel, leather & $-$0.235 & 0.015 & 0.260 & 0.275 & 0.011 \\
Utilities and sanitary services & 0.152 & 0.011 & 0.267 & 0.383 & 0.023 \\
Mining & 0.067 & 0.011 & 0.277 & 0.318 & 0.009 \\
Public administration (incl. postal) & 0.142 & 0.009 & 0.250 & 0.330 & 0.075 \\
Manufacturing: food, tobacco & $-$0.141 & 0.007 & 0.264 & 0.300 & 0.021 \\
Manufacturing: electrical equipment & 0.186 & 0.003 & 0.230 & 0.298 & 0.021 \\
Manufacturing: paper, printing, publishing & 0.034 & 0.002 & 0.229 & 0.333 & 0.026 \\
Computer and data processing services & 0.460 & 0.001 & 0.233 & 0.327 & 0.017 \\
Manufacturing: primary and fabricated metals & $-$0.061 & 0.001 & 0.222 & 0.269 & 0.029 \\
Manufacturing: wood, furniture, stone, glass & $-$0.183 & 0.000 & 0.227 & 0.289 & 0.025 \\
Manufacturing: machinery and computers & 0.105 & $-$0.002 & 0.235 & 0.294 & 0.033 \\
Social services, religious, membership & $-$0.128 & $-$0.006 & 0.260 & 0.236 & 0.019 \\
Wholesale trade & $-$0.008 & $-$0.008 & 0.230 & 0.289 & 0.053 \\
Construction & $-$0.031 & $-$0.010 & 0.205 & 0.282 & 0.096 \\
Transportation (excl. postal) & $-$0.052 & $-$0.019 & 0.195 & 0.278 & 0.058 \\
Eating, drinking, hotels & $-$0.467 & $-$0.023 & 0.207 & 0.311 & 0.039 \\
Retail trade (excl. eating, drinking) & $-$0.177 & $-$0.035 & 0.197 & 0.285 & 0.090 \\
Real estate & $-$0.072 & $-$0.036 & 0.168 & 0.177 & 0.014 \\
Automotive and repair services & $-$0.205 & $-$0.043 & 0.172 & 0.250 & 0.021 \\
Other business services & $-$0.173 & $-$0.043 & 0.183 & 0.228 & 0.029 \\
Personal services, entertainment & $-$0.195 & $-$0.046 & 0.180 & 0.237 & 0.020 \\
Agriculture, forestry, fishing & $-$0.530 & $-$0.064 & 0.100 & 0.177 & 0.019 \\
\bottomrule
\end{tabular}}}
\par\smallskip
\begin{minipage}{\textwidth}\footnotesize \textit{Notes:} Industries harmonised to the ind1990ddx scheme (222 industries, 33 groups). Industry level $P_j$ and career slope $B_j$ from CPS MORG men aged 25--64, 1992--2002 (743,983 observations): relative log wage $= f(\text{age}) + P_j + B_j(\text{age}-40)/10$, both employment-weighted mean zero. Rows sorted by $B$. Raw growth: mean relative log wage at 45--55 minus 25--34. Return to experience: predicted wage growth from 5 to 25 years of potential experience with years of schooling held fixed. Industry averages; the composition of workers is not held fixed. Across men, the standard deviation of $P$ is 0.232 and of $B$ 0.037. Source: \texttt{IGE/32.IndustryCareers.py} (\texttt{output/industry\_careers.csv}); table built by \texttt{IGE/28.PaperBTables.py}.\end{minipage}
\end{table}

\begin{table}[!htbp]
\centering
\caption{Sorting into industries by parental income (PSID, sons)}
\label{tab:industrysorting}
\resizebox{\ifdim\width>\textwidth\textwidth\else\width\fi}{!}{%
\resizebox*{!}{\ifdim\totalheight>0.85\textheight 0.85\textheight\else\totalheight\fi}{%
\begin{tabular}{lccccccc}
\toprule
Panel / row & \begin{tabular}[b]{@{}c@{}}Industry\\level $P$\end{tabular} & \begin{tabular}[b]{@{}c@{}}Industry\\slope $B$\end{tabular} & \begin{tabular}[b]{@{}c@{}}Finance,\\insurance,\\real estate\end{tabular} & \begin{tabular}[b]{@{}c@{}}Legal,\\professional\\services\end{tabular} & \begin{tabular}[b]{@{}c@{}}Computer,\\business\\services\end{tabular} & \begin{tabular}[b]{@{}c@{}}Steep-slope\\industry\end{tabular} & \begin{tabular}[b]{@{}c@{}}High-premium\\industry\end{tabular} \\
 & (1) & (2) & (3) & (4) & (5) & (6) & (7) \\
\midrule
\multicolumn{8}{l}{\textit{Pooled (weighted)}} \\
ln parental income, sons 25--34 & $0.077^{***}$ & $0.010^{***}$ & $0.036^{***}$ & $0.034^{***}$ & $0.022^{***}$ & $0.087^{***}$ & $0.085^{***}$ \\
 & (0.009) & (0.001) & (0.007) & (0.005) & (0.006) & (0.013) & (0.014) \\
\quad Mean of outcome, 25--34 & $-$0.007 & $-$0.002 & 0.049 & 0.038 & 0.052 & 0.257 & 0.239 \\
ln parental income, sons 35--55 & $0.075^{***}$ & $0.009^{***}$ & $0.037^{***}$ & $0.034^{***}$ & $0.023^{***}$ & $0.081^{***}$ & $0.082^{***}$ \\
 & (0.011) & (0.002) & (0.010) & (0.007) & (0.009) & (0.019) & (0.019) \\
\quad Mean of outcome, 35--55 & 0.021 & 0.001 & 0.054 & 0.041 & 0.057 & 0.293 & 0.265 \\
Change, (35--55) $-$ (25--34) & $-0.002$ & $-0.001$ & $0.001$ & $0.000$ & $0.001$ & $-0.006$ & $-0.003$ \\
 & (0.007) & (0.001) & (0.008) & (0.006) & (0.008) & (0.015) & (0.013) \\
\addlinespace
\multicolumn{8}{l}{\textit{Black}} \\
ln parental income, sons 25--34 & $0.044^{***}$ & $0.005^{***}$ & $0.017^{***}$ & $0.008^{***}$ & $0.008$ & $0.062^{***}$ & $0.052^{***}$ \\
 & (0.008) & (0.001) & (0.007) & (0.003) & (0.007) & (0.016) & (0.015) \\
\quad Mean of outcome, 25--34 & $-$0.062 & $-$0.007 & 0.039 & 0.010 & 0.048 & 0.222 & 0.199 \\
ln parental income, sons 35--55 & $0.042^{***}$ & $0.003^{*}$ & $0.015^{**}$ & $0.004$ & $0.011$ & $0.046^{**}$ & $0.060^{***}$ \\
 & (0.010) & (0.002) & (0.007) & (0.004) & (0.008) & (0.022) & (0.021) \\
\quad Mean of outcome, 35--55 & $-$0.042 & $-$0.006 & 0.036 & 0.013 & 0.045 & 0.227 & 0.215 \\
Change, (35--55) $-$ (25--34) & $-0.002$ & $-0.002$ & $-0.002$ & $-0.004$ & $0.003$ & $-0.016$ & $0.008$ \\
 & (0.009) & (0.002) & (0.006) & (0.005) & (0.009) & (0.019) & (0.018) \\
\addlinespace
\multicolumn{8}{l}{\textit{White}} \\
ln parental income, sons 25--34 & $0.082^{***}$ & $0.010^{***}$ & $0.041^{***}$ & $0.037^{***}$ & $0.030^{***}$ & $0.103^{***}$ & $0.101^{***}$ \\
 & (0.010) & (0.001) & (0.008) & (0.006) & (0.006) & (0.015) & (0.015) \\
\quad Mean of outcome, 25--34 & $-$0.009 & $-$0.002 & 0.048 & 0.039 & 0.048 & 0.249 & 0.235 \\
ln parental income, sons 35--55 & $0.084^{***}$ & $0.009^{***}$ & $0.048^{***}$ & $0.040^{***}$ & $0.036^{***}$ & $0.092^{***}$ & $0.100^{***}$ \\
 & (0.013) & (0.002) & (0.012) & (0.009) & (0.009) & (0.021) & (0.020) \\
\quad Mean of outcome, 35--55 & 0.022 & 0.001 & 0.053 & 0.043 & 0.053 & 0.292 & 0.265 \\
Change, (35--55) $-$ (25--34) & $0.003$ & $-0.001$ & $0.007$ & $0.003$ & $0.006$ & $-0.011$ & $-0.001$ \\
 & (0.008) & (0.001) & (0.009) & (0.007) & (0.008) & (0.017) & (0.015) \\
\addlinespace
\multicolumn{8}{l}{\textit{White $-$ Black}} \\
ln parental income, sons 25--34 & $0.038^{***}$ & $0.005^{***}$ & $0.024^{**}$ & $0.029^{***}$ & $0.022^{**}$ & $0.040^{*}$ & $0.048^{**}$ \\
 & (0.013) & (0.002) & (0.010) & (0.007) & (0.009) & (0.022) & (0.021) \\
ln parental income, sons 35--55 & $0.043^{**}$ & $0.006^{**}$ & $0.032^{**}$ & $0.036^{***}$ & $0.025^{**}$ & $0.045$ & $0.040$ \\
 & (0.017) & (0.003) & (0.014) & (0.009) & (0.012) & (0.030) & (0.029) \\
Change, (35--55) $-$ (25--34) & $0.005$ & $0.001$ & $0.008$ & $0.007$ & $0.003$ & $0.005$ & $-0.008$ \\
 & (0.012) & (0.002) & (0.011) & (0.008) & (0.012) & (0.026) & (0.023) \\
\bottomrule
\end{tabular}}}
\par\smallskip
\begin{minipage}{\textwidth}\footnotesize \textit{Notes:} Son-year regressions of the outcome in the header on log parental income with the controls of Table~\ref{tab:industrydecomp}; columns (3)--(7) are linear probability models. Industries harmonised to the ind1990ddx scheme (222 industries, 33 groups). Industry level $P_j$ and career slope $B_j$ from CPS MORG men aged 25--64, 1992--2002 (743,983 observations): relative log wage $= f(\text{age}) + P_j + B_j(\text{age}-40)/10$, both employment-weighted mean zero. Steep-slope and high-premium industries: top quartile of men's employment by $B_j$ and by $P_j$ (Table~\ref{tab:industrygroups}). Standard errors clustered by 1968 family. Pooled: weighted; Black and White: unweighted. $^{***}$ $p<0.01$, $^{**}$ $p<0.05$, $^{*}$ $p<0.1$. Source: \texttt{32.IndustryCareers.py} (\texttt{output/industry\_careers.csv}); table built by \texttt{28.PaperBTables.py}.\end{minipage}
\end{table}

\begin{table}[!htbp]
\centering
\caption{Wage persistence by age: industry level, industry career slope and within industry (PSID, sons)}
\label{tab:industrydecomp}
\resizebox{\ifdim\width>\textwidth\textwidth\else\width\fi}{!}{%
\resizebox*{!}{\ifdim\totalheight>0.85\textheight 0.85\textheight\else\totalheight\fi}{%
\begin{tabular}{lccccccccc}
\toprule
Panel / row & \begin{tabular}[b]{@{}c@{}}ln wage\end{tabular} & \begin{tabular}[b]{@{}c@{}}Industry\\level $P$\end{tabular} & \begin{tabular}[b]{@{}c@{}}Industry\\career slope\\$B\times t$\end{tabular} & \begin{tabular}[b]{@{}c@{}}Within-\\industry\\residual\end{tabular} & \begin{tabular}[b]{@{}c@{}}Occupation\\score\end{tabular} & \begin{tabular}[b]{@{}c@{}}Joint:\\occupation\end{tabular} & \begin{tabular}[b]{@{}c@{}}Joint:\\industry level\\net of occ.\end{tabular} & \begin{tabular}[b]{@{}c@{}}Joint:\\industry slope\\net of occ.\end{tabular} & \begin{tabular}[b]{@{}c@{}}Joint:\\residual\end{tabular} \\
 & (1) & (2) & (3) & (4) & (5) & (6) & (7) & (8) & (9) \\
\midrule
\multicolumn{10}{l}{\textit{Pooled (weighted): coefficient on log parental income}} \\
Sons aged 25--34 & $0.256^{***}$ & $0.077^{***}$ & $-0.010^{***}$ & $0.190^{***}$ & $0.159^{***}$ & $0.141^{***}$ & $0.025^{***}$ & $-0.011^{***}$ & $0.101^{***}$ \\
 & (0.017) & (0.009) & (0.001) & (0.015) & (0.010) & (0.008) & (0.004) & (0.001) & (0.014) \\
Sons aged 35--55 & $0.412^{***}$ & $0.075^{***}$ & $0.003^{***}$ & $0.335^{***}$ & $0.173^{***}$ & $0.155^{***}$ & $0.021^{***}$ & $0.003^{***}$ & $0.233^{***}$ \\
 & (0.028) & (0.011) & (0.001) & (0.024) & (0.013) & (0.010) & (0.005) & (0.001) & (0.022) \\
\addlinespace
\multicolumn{10}{l}{\textit{Black: coefficient on log parental income}} \\
Sons aged 25--34 & $0.148^{***}$ & $0.044^{***}$ & $-0.005^{***}$ & $0.109^{***}$ & $0.071^{***}$ & $0.060^{***}$ & $0.013^{***}$ & $-0.006^{***}$ & $0.081^{***}$ \\
 & (0.021) & (0.008) & (0.001) & (0.018) & (0.011) & (0.009) & (0.005) & (0.001) & (0.017) \\
Sons aged 35--55 & $0.157^{***}$ & $0.042^{***}$ & $-0.001$ & $0.116^{***}$ & $0.053^{***}$ & $0.044^{***}$ & $0.015^{***}$ & $0.000$ & $0.098^{***}$ \\
 & (0.031) & (0.010) & (0.001) & (0.027) & (0.014) & (0.012) & (0.005) & (0.001) & (0.023) \\
\addlinespace
\multicolumn{10}{l}{\textit{White: coefficient on log parental income}} \\
Sons aged 25--34 & $0.241^{***}$ & $0.082^{***}$ & $-0.011^{***}$ & $0.170^{***}$ & $0.157^{***}$ & $0.141^{***}$ & $0.024^{***}$ & $-0.013^{***}$ & $0.089^{***}$ \\
 & (0.020) & (0.010) & (0.001) & (0.017) & (0.012) & (0.010) & (0.005) & (0.001) & (0.016) \\
Sons aged 35--55 & $0.419^{***}$ & $0.084^{***}$ & $0.002^{***}$ & $0.332^{***}$ & $0.170^{***}$ & $0.154^{***}$ & $0.022^{***}$ & $0.002^{***}$ & $0.241^{***}$ \\
 & (0.034) & (0.013) & (0.001) & (0.028) & (0.015) & (0.012) & (0.006) & (0.001) & (0.025) \\
\addlinespace
\multicolumn{10}{l}{\textit{White $-$ Black: coefficient on log parental income}} \\
Sons aged 25--34 & $0.093^{***}$ & $0.038^{***}$ & $-0.006^{***}$ & $0.061^{**}$ & $0.087^{***}$ & $0.081^{***}$ & $0.010$ & $-0.007^{***}$ & $0.008$ \\
 & (0.029) & (0.013) & (0.002) & (0.025) & (0.016) & (0.013) & (0.007) & (0.002) & (0.023) \\
Sons aged 35--55 & $0.262^{***}$ & $0.043^{**}$ & $0.003^{***}$ & $0.217^{***}$ & $0.117^{***}$ & $0.110^{***}$ & $0.007$ & $0.003^{***}$ & $0.143^{***}$ \\
 & (0.046) & (0.017) & (0.001) & (0.039) & (0.020) & (0.017) & (0.008) & (0.001) & (0.035) \\
\addlinespace
\multicolumn{10}{l}{\textit{Change in the coefficient, (35--55) $-$ (25--34)}} \\
Pooled (weighted) & $0.156^{***}$ & $-0.002$ & $0.013^{***}$ & $0.145^{***}$ & $0.014$ & $0.015^{*}$ & $-0.004$ & $0.014^{***}$ & $0.132^{***}$ \\
 & (0.024) & (0.007) & (0.002) & (0.022) & (0.010) & (0.009) & (0.004) & (0.002) & (0.021) \\
Black & $0.009$ & $-0.002$ & $0.005^{***}$ & $0.006$ & $-0.018$ & $-0.016^{*}$ & $0.002$ & $0.006^{***}$ & $0.017$ \\
 & (0.026) & (0.009) & (0.002) & (0.025) & (0.011) & (0.010) & (0.005) & (0.002) & (0.022) \\
White & $0.178^{***}$ & $0.003$ & $0.013^{***}$ & $0.162^{***}$ & $0.012$ & $0.014$ & $-0.002$ & $0.015^{***}$ & $0.152^{***}$ \\
 & (0.027) & (0.008) & (0.002) & (0.025) & (0.012) & (0.010) & (0.004) & (0.002) & (0.023) \\
White $-$ Black & $0.170^{***}$ & $0.005$ & $0.009^{***}$ & $0.156^{***}$ & $0.030^{*}$ & $0.029^{**}$ & $-0.004$ & $0.009^{***}$ & $0.135^{***}$ \\
 & (0.038) & (0.012) & (0.002) & (0.035) & (0.017) & (0.014) & (0.007) & (0.002) & (0.032) \\
\addlinespace
\multicolumn{10}{l}{\textit{Share of the change in column (1)}} \\
Pooled change &  & $-0.010$ & $0.084^{***}$ & $0.926^{***}$ & $0.092$ & $0.095^{*}$ & $-0.026$ & $0.087^{***}$ & $0.844^{***}$ \\
 &  & (0.047) & (0.014) & (0.044) & (0.062) & (0.052) & (0.027) & (0.014) & (0.060) \\
White $-$ Black change &  & $0.030$ & $0.051^{***}$ & $0.919^{***}$ & $0.179^{**}$ & $0.173^{**}$ & $-0.021$ & $0.054^{***}$ & $0.794^{***}$ \\
 &  & (0.071) & (0.016) & (0.070) & (0.087) & (0.074) & (0.041) & (0.016) & (0.084) \\
\bottomrule
\end{tabular}}}
\par\smallskip
\begin{minipage}{\textwidth}\footnotesize \textit{Notes:} One regression per cell: the son's outcome (column header) on log parental income, son-years in the age band, controls for quartics in the father's age at 16 and the son's age. Industries harmonised to the ind1990ddx scheme (222 industries, 33 groups). Industry level $P_j$ and career slope $B_j$ from CPS MORG men aged 25--64, 1992--2002 (743,983 observations): relative log wage $= f(\text{age}) + P_j + B_j(\text{age}-40)/10$, both employment-weighted mean zero. Log wage $=P_j+B_j t+$ within-industry residual, so columns (2)--(4) add up to column (1); column (5) is the occupation score of Table~\ref{tab:careerdecomp} on the same son-years; columns (6)--(9) split the log wage with occupation fixed effects and industry levels and slopes estimated jointly, so (7)$+$(8) is what industry adds beyond occupation. Same sample in every column: 1,822 Black sons (12,576 son-years) and 2,639 White sons (25,248). All cells are one stacked system (standard errors clustered by 1968 family); shares by the delta method. Pooled: weighted; Black and White: unweighted. Sign convention: White $-$ Black. Sons select into industries, so the split is accounting, not causal. $^{***}$ $p<0.01$, $^{**}$ $p<0.05$, $^{*}$ $p<0.1$. Source: \texttt{32.IndustryCareers.py} (\texttt{output/industry\_careers.csv}); table built by \texttt{28.PaperBTables.py}.\end{minipage}
\end{table}

\begin{table}[!htbp]
\centering
\caption{Industry and occupation $\times$ industry career profiles: linear and nonlinear (PSID, sons)}
\label{tab:industrynonlinear}
\resizebox{\ifdim\width>\textwidth\textwidth\else\width\fi}{!}{%
\resizebox*{!}{\ifdim\totalheight>0.85\textheight 0.85\textheight\else\totalheight\fi}{%
\begin{tabular}{lcccc}
\toprule
Piece of the log wage & \begin{tabular}[b]{@{}c@{}}Pooled\\change\end{tabular} & \begin{tabular}[b]{@{}c@{}}Share of\\pooled rise\end{tabular} & \begin{tabular}[b]{@{}c@{}}Change of\\White $-$ Black\\gap\end{tabular} & \begin{tabular}[b]{@{}c@{}}Share of\\gap change\end{tabular} \\
 & (1) & (2) & (3) & (4) \\
\midrule
\multicolumn{5}{l}{\textit{A. Sample of Table~\ref{tab:industrydecomp}}} \\
Log wage (total) & $0.155^{***}$ &  & $0.170^{***}$ &  \\
 & (0.024) &  & (0.038) &  \\
Current industry: linear slope $B_j t$ & $0.013^{***}$ & $0.084^{***}$ & $0.009^{***}$ & $0.050^{***}$ \\
 & (0.002) & (0.014) & (0.002) & (0.016) \\
Current industry: age-band profile & $0.018^{***}$ & $0.118^{***}$ & $0.013^{***}$ & $0.074^{***}$ \\
 & (0.002) & (0.020) & (0.003) & (0.023) \\
Current industry: piecewise profile (knots 30/35/45) & $0.018^{***}$ & $0.117^{***}$ & $0.012^{***}$ & $0.072^{***}$ \\
 & (0.002) & (0.019) & (0.003) & (0.023) \\
\addlinespace
\multicolumn{5}{l}{\textit{B. Sons with an industry and cell at 25--34}} \\
Log wage (total) & $0.146^{***}$ &  & $0.159^{***}$ &  \\
 & (0.024) &  & (0.039) &  \\
Industry at 25--34: linear slope & $0.014^{***}$ & $0.093^{***}$ & $0.009^{***}$ & $0.058^{***}$ \\
 & (0.002) & (0.017) & (0.003) & (0.022) \\
Industry at 25--34: age-band profile & $0.019^{***}$ & $0.132^{***}$ & $0.014^{***}$ & $0.086^{***}$ \\
 & (0.003) & (0.024) & (0.004) & (0.032) \\
Current occupation $\times$ industry cell: age-band profile & $0.036^{***}$ & $0.246^{***}$ & $0.022^{***}$ & $0.138^{***}$ \\
 & (0.003) & (0.038) & (0.004) & (0.035) \\
Current occupation group only: age-band profile & $0.031^{***}$ & $0.215^{***}$ & $0.021^{***}$ & $0.131^{***}$ \\
 & (0.002) & (0.034) & (0.003) & (0.032) \\
\quad What industry adds within the current cell & $0.005^{***}$ & $0.031^{***}$ & $0.001$ & $0.007$ \\
 & (0.001) & (0.010) & (0.002) & (0.015) \\
Cell at 25--34: age-band profile & $0.039^{***}$ & $0.267^{***}$ & $0.024^{***}$ & $0.148^{***}$ \\
 & (0.003) & (0.043) & (0.005) & (0.042) \\
Occupation group at 25--34 only: age-band profile & $0.033^{***}$ & $0.222^{***}$ & $0.019^{***}$ & $0.120^{***}$ \\
 & (0.003) & (0.036) & (0.004) & (0.033) \\
\quad What industry adds within the cell at 25--34 & $0.007^{***}$ & $0.045^{***}$ & $0.004$ & $0.028$ \\
 & (0.002) & (0.014) & (0.003) & (0.019) \\
\bottomrule
\end{tabular}}}
\par\smallskip
\begin{minipage}{\textwidth}\footnotesize \textit{Notes:} Each row is a piece of the son's log wage regressed on log parental income by age band (stacked, controls as in Table~\ref{tab:industrydecomp}); columns (1) and (3) give the change of the coefficient from 25--34 to 35--55, pooled and for the White $-$ Black gap, and columns (2) and (4) its share of the change for the log wage. Linear: $B_j t$ as in Table~\ref{tab:industrydecomp}. Age-band profile: the CPS MORG relative wage of the son's industry or cell in his five-year age band, around its age-weighted level (men 25--64, 1992--2002). Cells: 11 occupation groups $\times$ 33 industry groups (267 cells with at least 200 men; smaller cells take their occupation group's profile). `At 25--34': the son's modal industry or cell at 25--34, evaluated at his current age. Standard errors clustered by 1968 family; shares by the delta method. Pooled: weighted. $^{***}$ $p<0.01$, $^{**}$ $p<0.05$, $^{*}$ $p<0.1$. Source: \texttt{IGE/35.EarlyIndustryLatePayoff.py} (\texttt{output/early\_industry\_late\_payoff.csv}); table built by \texttt{IGE/28.PaperBTables.py}.\end{minipage}
\end{table}

\begin{table}[!htbp]
\centering
\caption{Wage growth with the same employer by parental income, within industries (PSID, sons)}
\label{tab:industrygrowth}
\resizebox{\ifdim\width>\textwidth\textwidth\else\width\fi}{!}{%
\resizebox*{!}{\ifdim\totalheight>0.85\textheight 0.85\textheight\else\totalheight\fi}{%
\begin{tabular}{lcccccc}
\toprule
Panel / row & \begin{tabular}[b]{@{}c@{}}All\\industries\end{tabular} & \begin{tabular}[b]{@{}c@{}}+ industry-\\group FE\end{tabular} & \begin{tabular}[b]{@{}c@{}}+ fine-\\industry FE\end{tabular} & \begin{tabular}[b]{@{}c@{}}Finance, legal,\\professional,\\business\end{tabular} & \begin{tabular}[b]{@{}c@{}}Other\\steep-slope\\industries\end{tabular} & \begin{tabular}[b]{@{}c@{}}All other\\industries\end{tabular} \\
 & (1) & (2) & (3) & (4) & (5) & (6) \\
\midrule
\multicolumn{7}{l}{\textit{Pooled (weighted)}} \\
ln parental income, sons 25--34 & $0.012^{**}$ & $0.007$ & $0.005$ & $0.024$ & $0.021$ & $0.000$ \\
 & (0.005) & (0.006) & (0.006) & (0.021) & (0.014) & (0.006) \\
\quad Stayer pairs, 25--34 & 9,365 & 9,365 & 9,365 & 1,121 & 1,834 & 6,410 \\
ln parental income, sons 35--55 & $0.001$ & $-0.002$ & $-0.002$ & $-0.001$ & $0.000$ & $-0.002$ \\
 & (0.003) & (0.003) & (0.003) & (0.010) & (0.005) & (0.004) \\
\quad Stayer pairs, 35--55 & 12,246 & 12,246 & 12,246 & 1,416 & 2,869 & 7,961 \\
\addlinespace
\multicolumn{7}{l}{\textit{Black}} \\
ln parental income, sons 25--34 & $0.003$ & $0.002$ & $0.002$ &  & $0.018$ & $0.001$ \\
 & (0.008) & (0.008) & (0.009) &  & (0.020) & (0.010) \\
\quad Stayer pairs, 25--34 & 2,947 & 2,947 & 2,947 &  & 610 & 2,100 \\
ln parental income, sons 35--55 & $0.000$ & $0.001$ & $0.002$ & $0.030$ & $0.018^{*}$ & $-0.009$ \\
 & (0.004) & (0.004) & (0.005) & (0.024) & (0.010) & (0.006) \\
\quad Stayer pairs, 35--55 & 3,864 & 3,864 & 3,864 & 309 & 857 & 2,698 \\
\addlinespace
\multicolumn{7}{l}{\textit{White}} \\
ln parental income, sons 25--34 & $0.021^{***}$ & $0.015^{***}$ & $0.015^{***}$ & $0.040^{**}$ & $0.026$ & $0.009$ \\
 & (0.006) & (0.006) & (0.006) & (0.017) & (0.019) & (0.006) \\
\quad Stayer pairs, 25--34 & 6,427 & 6,427 & 6,427 & 884 & 1,224 & 4,319 \\
ln parental income, sons 35--55 & $0.001$ & $-0.002$ & $-0.003$ & $0.004$ & $-0.008$ & $0.000$ \\
 & (0.004) & (0.004) & (0.004) & (0.012) & (0.007) & (0.005) \\
\quad Stayer pairs, 35--55 & 8,383 & 8,383 & 8,383 & 1,107 & 2,012 & 5,264 \\
\addlinespace
\multicolumn{7}{l}{\textit{White $-$ Black}} \\
ln parental income, sons 25--34 & $0.018^{*}$ & $0.013$ & $0.014$ &  & $0.008$ & $0.008$ \\
 & (0.010) & (0.010) & (0.011) &  & (0.027) & (0.011) \\
ln parental income, sons 35--55 & $0.001$ & $-0.003$ & $-0.005$ & $-0.026$ & $-0.025^{**}$ & $0.009$ \\
 & (0.006) & (0.006) & (0.007) & (0.027) & (0.012) & (0.008) \\
\bottomrule
\end{tabular}}}
\par\smallskip
\begin{minipage}{\textwidth}\footnotesize \textit{Notes:} Outcome in every column: annualised log-wage growth between consecutive interviews 1--2 years apart (waves 1985--2023) for sons who stayed with their employer. Columns (2)--(3) add fixed effects for the industry of the later interview, so they compare sons in the same industry; columns (4)--(6) split the sample by the son's industry. Controls: quartic in age, year fixed effects. Cells with fewer than 250 pairs are left blank. Standard errors clustered by 1968 family. $^{***}$ $p<0.01$, $^{**}$ $p<0.05$, $^{*}$ $p<0.1$. Source: \texttt{32.IndustryCareers.py} (\texttt{output/industry\_careers.csv}); table built by \texttt{28.PaperBTables.py}.\end{minipage}
\end{table}

\begin{table}[!htbp]
\centering
\caption{Following the father into his industry: how often and does it pay (PSID, sons)}
\label{tab:fatherindustry}
\resizebox{\ifdim\width>\textwidth\textwidth\else\width\fi}{!}{%
\resizebox*{!}{\ifdim\totalheight>0.85\textheight 0.85\textheight\else\totalheight\fi}{%
\begin{tabular}{lccccccc}
\toprule
Panel / row & \begin{tabular}[b]{@{}c@{}}Son in\\finance,\\insurance,\\real estate\end{tabular} & \begin{tabular}[b]{@{}c@{}}Son in\\legal,\\professional\end{tabular} & \begin{tabular}[b]{@{}c@{}}ln wage,\\25--34\end{tabular} & \begin{tabular}[b]{@{}c@{}}ln wage,\\35--55\end{tabular} & \begin{tabular}[b]{@{}c@{}}Annual\\wage growth,\\30--55\end{tabular} & \begin{tabular}[b]{@{}c@{}}Same-employer\\wage growth,\\30--55\end{tabular} & \begin{tabular}[b]{@{}c@{}}ln wage,\\person FE\end{tabular} \\
 & (1) & (2) & (3) & (4) & (5) & (6) & (7) \\
\midrule
\multicolumn{8}{l}{\textit{Pooled (weighted)}} \\
Father worked in the same industry set & $0.108^{**}$ & $0.072^{*}$ &  &  &  &  &  \\
 & (0.047) & (0.040) &  &  &  &  &  \\
Son in the father's industry group &  &  & $0.079^{***}$ & $0.054^{*}$ & $0.005$ & $-0.006$ & $0.039$ \\
 &  &  & (0.023) & (0.030) & (0.005) & (0.005) & (0.025) \\
Tenure (per 10 years) $\times$ in father's industry &  &  &  &  &  &  & $-0.042^{**}$ \\
 &  &  &  &  &  &  & (0.020) \\
\addlinespace
\multicolumn{8}{l}{\textit{Black}} \\
Father worked in the same industry set & $-0.003$ & $-0.012^{***}$ &  &  &  &  &  \\
 & (0.023) & (0.004) &  &  &  &  &  \\
Son in the father's industry group &  &  & $0.034$ & $0.088$ & $0.000$ & $-0.021^{**}$ & $0.115^{***}$ \\
 &  &  & (0.042) & (0.073) & (0.010) & (0.010) & (0.040) \\
Tenure (per 10 years) $\times$ in father's industry &  &  &  &  &  &  & $-0.068^{*}$ \\
 &  &  &  &  &  &  & (0.039) \\
\addlinespace
\multicolumn{8}{l}{\textit{White}} \\
Father worked in the same industry set & $0.122^{***}$ & $0.067^{*}$ &  &  &  &  &  \\
 & (0.047) & (0.040) &  &  &  &  &  \\
Son in the father's industry group &  &  & $0.057^{**}$ & $0.046$ & $0.006$ & $-0.004$ & $0.024$ \\
 &  &  & (0.023) & (0.029) & (0.006) & (0.005) & (0.022) \\
Tenure (per 10 years) $\times$ in father's industry &  &  &  &  &  &  & $-0.040^{**}$ \\
 &  &  &  &  &  &  & (0.019) \\
\addlinespace
\multicolumn{8}{l}{\textit{White $-$ Black}} \\
Father worked in the same industry set & $0.125^{**}$ & $0.079^{*}$ &  &  &  &  &  \\
 & (0.052) & (0.041) &  &  &  &  &  \\
Son in the father's industry group &  &  & $0.023$ & $-0.041$ & $0.006$ & $0.016$ &  \\
 &  &  & (0.048) & (0.078) & (0.012) & (0.012) &  \\
Tenure (per 10 years) $\times$ in father's industry &  &  &  &  &  &  & $0.028$ \\
 &  &  &  &  &  &  & (0.043) \\
\bottomrule
\end{tabular}}}
\par\smallskip
\begin{minipage}{\textwidth}\footnotesize \textit{Notes:} Father: male head of the son's family at 16; his industry is the modal code over his years as head when the son was 12--18 (father-present sons: 52\% of Black and 87\% of White sons). The PSID has no employer identifiers, so following is measured at the industry level. Columns (1)--(2): son-years 25--55, linear probability that the son works in the industry set, on an indicator that the father did and log parental income. Columns (3)--(6): son-year regressions on an indicator that the son works in his father's industry group (33 groups), holding log parental income, years of education and the level $P_j$ and slope $B_j$ of the son's industry fixed; growth between consecutive interviews 1--2 years apart. Column (7): within-son regression on tenure, its square, tenure $\times$ following, following and tenure $\times$ log parental income, with a quartic in age and year fixed effects. Standard errors clustered by 1968 family. Pooled: weighted. Only 21 Black sons had a father in finance, insurance or real estate and 8 in legal or professional services. Sons who follow may differ in unobserved ways; descriptive. $^{***}$ $p<0.01$, $^{**}$ $p<0.05$, $^{*}$ $p<0.1$. Source: \texttt{32.IndustryCareers.py} (\texttt{output/industry\_careers.csv}); table built by \texttt{28.PaperBTables.py}.\end{minipage}
\end{table}


\subsection{Growth, tenure, promotion and displacement}

Table~\ref{tab:careergrowth} splits annual wage growth into occupational-score and residual growth with the same and a
new employer; Table~\ref{tab:tenure} reports wage--tenure profiles within sons, which steepen with parental income for
White sons (0.068$^{***}$ per ten years of tenure and log point of parental income) and not for Black sons (0.013).
Tables~\ref{tab:promotion} and \ref{tab:hegde} report entry into supervisory and managerial occupations and the
analogue to \citet{hegde2026racial}: the after-35 change relative to 25--34 is not significant (triple interaction
0.008, s.e.\ 0.011, for the change of the White $-$ Black difference), and holding the wage and occupation at 28--32 fixed the gradients are 0.014 for White and
$-0.010$ for Black sons. Table~\ref{tab:careerdisp} reports the wage loss after displacement, 0.131$^{***}$, of which
0.028$^{***}$ is a lower occupational score. Table~\ref{tab:mediation} reports the coefficient movements when education
and instability shares are added to the IGE, and Table~\ref{tab:nlsycareer} the NLSY79 version of the occupational
decomposition: the pooled wage gradient rises from 0.252$^{***}$ at 25--34 to 0.381$^{***}$ at 45--55, almost all within
occupations, and the White--Black gap changes by 0.059$^{*}$.

\begin{table}[!htbp]
\centering
\caption{Annual wage growth by parental income: within and across employers (PSID, sons)}
\label{tab:careergrowth}
\resizebox{\ifdim\width>\textwidth\textwidth\else\width\fi}{!}{%
\resizebox*{!}{\ifdim\totalheight>0.85\textheight 0.85\textheight\else\totalheight\fi}{%
\begin{tabular}{lcccccc}
\toprule
Panel / row & \begin{tabular}[b]{@{}c@{}}Annual\\ln wage\\growth\end{tabular} & \begin{tabular}[b]{@{}c@{}}Score growth,\\same\\employer\end{tabular} & \begin{tabular}[b]{@{}c@{}}Residual growth,\\same\\employer\end{tabular} & \begin{tabular}[b]{@{}c@{}}Score growth,\\new\\employer\end{tabular} & \begin{tabular}[b]{@{}c@{}}Residual growth,\\new\\employer\end{tabular} & \begin{tabular}[b]{@{}c@{}}Stayed\\with\\employer\end{tabular} \\
 & (1) & (2) & (3) & (4) & (5) & (6) \\
\midrule
\multicolumn{7}{l}{\textit{Pooled (weighted)}} \\
ln parental income, sons 25--34 & $0.015^{***}$ & $0.001$ & $0.010^{**}$ & $0.002$ & $0.003$ & $0.022^{**}$ \\
 & (0.005) & (0.002) & (0.004) & (0.001) & (0.003) & (0.009) \\
\quad Mean of outcome / pairs, 25--34 & 0.054 & 0.009 & 0.039 & 0.002 & 0.004 & 11,578 pairs \\
ln parental income, sons 35--55 & $0.004$ & $-0.002^{*}$ & $0.003$ & $0.000$ & $0.002$ & $0.020^{**}$ \\
 & (0.003) & (0.001) & (0.003) & (0.001) & (0.002) & (0.008) \\
\quad Mean of outcome / pairs, 35--55 & 0.020 & 0.002 & 0.017 & $-$0.001 & 0.002 & 14,185 pairs \\
\addlinespace
\multicolumn{7}{l}{\textit{Black}} \\
ln parental income, sons 25--34 & $-0.001$ & $0.004$ & $0.000$ & $0.003$ & $-0.008$ & $0.018$ \\
 & (0.006) & (0.003) & (0.007) & (0.002) & (0.005) & (0.013) \\
\quad Mean of outcome / pairs, 25--34 & 0.042 & 0.003 & 0.040 & 0.001 & $-$0.003 & 3,725 pairs \\
ln parental income, sons 35--55 & $-0.002$ & $-0.001$ & $0.001$ & $-0.001$ & $-0.001$ & $-0.011$ \\
 & (0.004) & (0.002) & (0.004) & (0.001) & (0.003) & (0.012) \\
\quad Mean of outcome / pairs, 35--55 & 0.015 & 0.003 & 0.013 & 0.000 & $-$0.002 & 4,537 pairs \\
\addlinespace
\multicolumn{7}{l}{\textit{White}} \\
ln parental income, sons 25--34 & $0.020^{***}$ & $0.000$ & $0.018^{***}$ & $0.001$ & $0.001$ & $0.020^{**}$ \\
 & (0.005) & (0.002) & (0.005) & (0.001) & (0.003) & (0.010) \\
\quad Mean of outcome / pairs, 25--34 & 0.054 & 0.010 & 0.038 & 0.002 & 0.005 & 7,862 pairs \\
ln parental income, sons 35--55 & $0.007^{**}$ & $-0.001$ & $0.003$ & $0.001$ & $0.004^{**}$ & $0.025^{***}$ \\
 & (0.003) & (0.001) & (0.004) & (0.001) & (0.002) & (0.009) \\
\quad Mean of outcome / pairs, 35--55 & 0.021 & 0.002 & 0.018 & $-$0.001 & 0.002 & 9,649 pairs \\
\addlinespace
\multicolumn{7}{l}{\textit{White $-$ Black}} \\
ln parental income, sons 25--34 & $0.021^{***}$ & $-0.003$ & $0.018^{**}$ & $-0.002$ & $0.009$ & $0.001$ \\
 & (0.008) & (0.004) & (0.008) & (0.003) & (0.006) & (0.016) \\
ln parental income, sons 35--55 & $0.009^{*}$ & $0.000$ & $0.002$ & $0.002$ & $0.005$ & $0.036^{**}$ \\
 & (0.005) & (0.002) & (0.005) & (0.001) & (0.004) & (0.015) \\
\bottomrule
\end{tabular}}}
\par\smallskip
\begin{minipage}{\textwidth}\footnotesize \textit{Notes:} Consecutive interviews 1--2 years apart (waves 1985--2023) with a valid wage, a mapped occupation and tenure with the present employer. Annualised log-wage growth = score growth + residual growth, each split by whether the son stayed with the same employer; columns (2)--(5) add up to column (1). Column (6): share of stayers (linear probability). Controls: quartic in age and year fixed effects. Mean row: mean of the outcome in columns (1)--(5) and the number of son-year pairs. Coding noise in occupations inflates score changes in both directions. Standard errors clustered by 1968 family. Occupational wage score: mean log hourly wage of men aged 25--64 in the occupation relative to all men (CPS MORG, base 1992--2002), occupations harmonised to Dorn's occ1990dd scheme; log wage = score (between occupations) + residual (within occupation), so the score and residual columns add up to the log-wage column on the same sample. $^{***}$ $p<0.01$, $^{**}$ $p<0.05$, $^{*}$ $p<0.1$. Source: \texttt{IGE/27.CareerWageGrowth.py} (\texttt{output/career\_wage\_growth.csv}); table built by \texttt{IGE/28.PaperBTables.py}.\end{minipage}
\end{table}

\begin{table}[!htbp]
\centering
\caption{Wage--tenure profiles by parental income (PSID, sons, person fixed effects)}
\label{tab:tenure}
\resizebox{\ifdim\width>\textwidth\textwidth\else\width\fi}{!}{%
\resizebox*{!}{\ifdim\totalheight>0.85\textheight 0.85\textheight\else\totalheight\fi}{%
\begin{tabular}{lccc}
\toprule
Panel / regressor & \begin{tabular}[b]{@{}c@{}}ln wage\end{tabular} & \begin{tabular}[b]{@{}c@{}}Occupation\\score\end{tabular} & \begin{tabular}[b]{@{}c@{}}Within-\\occupation\\residual\end{tabular} \\
 & (1) & (2) & (3) \\
\midrule
\multicolumn{4}{l}{\textit{Pooled (weighted)}} \\
Tenure (per 10 years) & $0.242^{***}$ & $0.037^{***}$ & $0.206^{***}$ \\
 & (0.019) & (0.009) & (0.021) \\
Tenure squared & $-0.076^{***}$ & $-0.007^{**}$ & $-0.068^{***}$ \\
 & (0.007) & (0.003) & (0.008) \\
Tenure $\times$ ln parental income & $0.032^{**}$ & $-0.015^{**}$ & $0.047^{***}$ \\
 & (0.014) & (0.008) & (0.016) \\
Son-years & 30,561 & 30,561 & 30,561 \\
\addlinespace
\multicolumn{4}{l}{\textit{Black}} \\
Tenure (per 10 years) & $0.268^{***}$ & $0.056^{***}$ & $0.212^{***}$ \\
 & (0.030) & (0.015) & (0.032) \\
Tenure squared & $-0.065^{***}$ & $-0.012^{**}$ & $-0.053^{***}$ \\
 & (0.010) & (0.005) & (0.011) \\
Tenure $\times$ ln parental income & $0.003$ & $-0.010$ & $0.013$ \\
 & (0.019) & (0.009) & (0.019) \\
Son-years & 10,217 & 10,217 & 10,217 \\
\addlinespace
\multicolumn{4}{l}{\textit{White}} \\
Tenure (per 10 years) & $0.233^{***}$ & $0.037^{***}$ & $0.196^{***}$ \\
 & (0.019) & (0.010) & (0.021) \\
Tenure squared & $-0.075^{***}$ & $-0.006$ & $-0.069^{***}$ \\
 & (0.007) & (0.004) & (0.008) \\
Tenure $\times$ ln parental income & $0.046^{***}$ & $-0.023^{***}$ & $0.068^{***}$ \\
 & (0.016) & (0.009) & (0.019) \\
Son-years & 20,357 & 20,357 & 20,357 \\
\addlinespace
\multicolumn{4}{l}{\textit{White $-$ Black}} \\
Tenure (per 10 years) & $-0.035$ & $-0.019$ & $-0.016$ \\
 & (0.035) & (0.018) & (0.038) \\
Tenure $\times$ ln parental income & $0.042^{*}$ & $-0.013$ & $0.055^{**}$ \\
 & (0.025) & (0.012) & (0.027) \\
\bottomrule
\end{tabular}}}
\par\smallskip
\begin{minipage}{\textwidth}\footnotesize \textit{Notes:} Within-son regressions (person fixed effects) of the outcome in the header on tenure with the present employer (per 10 years), its square and tenure interacted with centred log parental income, with a quartic in age and year fixed effects (waves 1984--2023). Within-son tenure slopes mix the return to seniority with the selection of good matches, so they are descriptive wage--tenure profiles. Standard errors clustered by 1968 family. Occupational wage score: mean log hourly wage of men aged 25--64 in the occupation relative to all men (CPS MORG, base 1992--2002), occupations harmonised to Dorn's occ1990dd scheme; log wage = score (between occupations) + residual (within occupation), so the score and residual columns add up to the log-wage column on the same sample. $^{***}$ $p<0.01$, $^{**}$ $p<0.05$, $^{*}$ $p<0.1$. Source: \texttt{27.CareerWageGrowth.py} (\texttt{output/career\_wage\_growth.csv}); table built by \texttt{28.PaperBTables.py}.\end{minipage}
\end{table}

\begin{table}[!htbp]
\centering
\caption{Promotion and occupational upgrading by parental income (PSID, sons)}
\label{tab:promotion}
\resizebox{\ifdim\width>\textwidth\textwidth\else\width\fi}{!}{%
\resizebox*{!}{\ifdim\totalheight>0.85\textheight 0.85\textheight\else\totalheight\fi}{%
\begin{tabular}{lcccccc}
\toprule
Panel / row & \begin{tabular}[b]{@{}c@{}}Moved to a\\higher-paying\\occupation\end{tabular} & \begin{tabular}[b]{@{}c@{}}Entered a\\supervisory or\\managerial occ.\end{tabular} & \begin{tabular}[b]{@{}c@{}}Annual\\change in\\score\end{tabular} & \begin{tabular}[b]{@{}c@{}}Reported\\promotion\\(1976--87)\end{tabular} & \begin{tabular}[b]{@{}c@{}}Position change,\\same employer\\(1984--2001)\end{tabular} & \begin{tabular}[b]{@{}c@{}}Position change\\with a\\wage gain\end{tabular} \\
 & (1) & (2) & (3) & (4) & (5) & (6) \\
\midrule
\multicolumn{7}{l}{\textit{Pooled (weighted)}} \\
ln parental income, sons 25--34 & $-0.031^{***}$ & $0.028^{***}$ & $0.002$ & $0.012^{*}$ & $0.012$ & $0.004$ \\
 & (0.007) & (0.005) & (0.002) & (0.006) & (0.009) & (0.007) \\
\quad Mean of outcome, 25--34 & 0.250 & 0.087 & 0.012 & 0.054 & 0.141 & 0.090 \\
ln parental income, sons 35--55 & $-0.024^{***}$ & $0.029^{***}$ & $-0.002^{*}$ &  & $0.033^{***}$ & $0.019^{**}$ \\
 & (0.006) & (0.006) & (0.001) &  & (0.011) & (0.008) \\
\quad Mean of outcome, 35--55 & 0.188 & 0.089 & 0.000 &  & 0.128 & 0.072 \\
\addlinespace
\multicolumn{7}{l}{\textit{Black}} \\
ln parental income, sons 25--34 & $-0.011$ & $0.012^{**}$ & $0.004^{*}$ & $0.009$ & $0.041^{***}$ & $0.021^{**}$ \\
 & (0.010) & (0.006) & (0.002) & (0.010) & (0.009) & (0.009) \\
\quad Mean of outcome, 25--34 & 0.269 & 0.043 & 0.006 & 0.035 & 0.104 & 0.061 \\
ln parental income, sons 35--55 & $-0.011$ & $0.000$ & $-0.002$ &  & $0.061^{***}$ & $0.022^{**}$ \\
 & (0.009) & (0.006) & (0.001) &  & (0.018) & (0.010) \\
\quad Mean of outcome, 35--55 & 0.214 & 0.053 & 0.003 &  & 0.102 & 0.061 \\
\addlinespace
\multicolumn{7}{l}{\textit{White}} \\
ln parental income, sons 25--34 & $-0.024^{***}$ & $0.024^{***}$ & $0.001$ & $0.015^{*}$ & $0.007$ & $0.003$ \\
 & (0.007) & (0.006) & (0.002) & (0.008) & (0.010) & (0.008) \\
\quad Mean of outcome, 25--34 & 0.242 & 0.089 & 0.012 & 0.055 & 0.140 & 0.088 \\
ln parental income, sons 35--55 & $-0.019^{**}$ & $0.025^{***}$ & $-0.001$ &  & $0.042^{***}$ & $0.026^{***}$ \\
 & (0.007) & (0.007) & (0.001) &  & (0.013) & (0.009) \\
\quad Mean of outcome, 35--55 & 0.187 & 0.095 & 0.000 &  & 0.125 & 0.069 \\
\addlinespace
\multicolumn{7}{l}{\textit{White $-$ Black}} \\
ln parental income, sons 25--34 & $-0.013$ & $0.012$ & $-0.004$ & $0.005$ & $-0.034^{**}$ & $-0.018$ \\
 & (0.012) & (0.008) & (0.003) & (0.012) & (0.014) & (0.012) \\
ln parental income, sons 35--55 & $-0.008$ & $0.024^{***}$ & $0.001$ &  & $-0.019$ & $0.004$ \\
 & (0.011) & (0.009) & (0.002) &  & (0.022) & (0.013) \\
\bottomrule
\end{tabular}}}
\par\smallskip
\begin{minipage}{\textwidth}\footnotesize \textit{Notes:} Son-year linear probability models (outcome in the header), male heads, controls: quartic in age and year fixed effects. Columns (1)--(3): change of the harmonised occupation between consecutive interviews; column (2) among sons not in a supervisory or managerial occupation at the previous interview. Column (4): the employed head whose job or position began within about a year reports that the previous one ended by promotion (waves 1976--1987, when most sons are under 40, hence no 35--55 estimate). Columns (5)--(6): the present position with the present employer began since the previous interview (self-reported, waves 1984--2001). Coded occupation changes include coding noise. Standard errors clustered by 1968 family. $^{***}$ $p<0.01$, $^{**}$ $p<0.05$, $^{*}$ $p<0.1$. Source: \texttt{IGE/27.CareerWageGrowth.py} (\texttt{output/career\_wage\_growth.csv}); table built by \texttt{IGE/28.PaperBTables.py}.\end{minipage}
\end{table}

\begin{table}[!htbp]
\centering
\caption{Do Black sons from higher-income families progress less after 35? (PSID, sons)}
\label{tab:hegde}
\resizebox{\ifdim\width>\textwidth\textwidth\else\width\fi}{!}{%
\resizebox*{!}{\ifdim\totalheight>0.85\textheight 0.85\textheight\else\totalheight\fi}{%
\begin{tabular}{lcccccc}
\toprule
Panel / regressor & \begin{tabular}[b]{@{}c@{}}ln wage\end{tabular} & \begin{tabular}[b]{@{}c@{}}Occupation\\score\end{tabular} & \begin{tabular}[b]{@{}c@{}}Within-\\occupation\\residual\end{tabular} & \begin{tabular}[b]{@{}c@{}}Entered\\supervisory or\\managerial\end{tabular} & \begin{tabular}[b]{@{}c@{}}Moved to a\\higher-paying\\occupation\end{tabular} & \begin{tabular}[b]{@{}c@{}}Position change,\\same\\employer\end{tabular} \\
 & (1) & (2) & (3) & (4) & (5) & (6) \\
\midrule
\multicolumn{7}{l}{\textit{A. Unconditional}} \\
ln parental income (centred) & $0.279^{***}$ & $0.162^{***}$ & $0.109^{***}$ & $0.024^{***}$ & $-0.023^{***}$ & $0.007$ \\
 & (0.019) & (0.011) & (0.016) & (0.006) & (0.007) & (0.010) \\
\quad $\times$ age 35--55 & $0.137^{***}$ & $0.013$ & $0.136^{***}$ & $-0.001$ & $0.003$ & $0.032^{**}$ \\
 & (0.026) & (0.012) & (0.023) & (0.008) & (0.010) & (0.015) \\
Black & $-0.166^{***}$ & $-0.121^{***}$ & $-0.039^{**}$ & $-0.031^{***}$ & $0.024^{***}$ & $-0.014$ \\
 & (0.020) & (0.012) & (0.017) & (0.006) & (0.008) & (0.010) \\
Black $\times$ age 35--55 & $-0.064^{***}$ & $-0.026^{**}$ & $-0.035^{*}$ & $-0.001$ & $-0.006$ & $0.037^{**}$ \\
 & (0.023) & (0.012) & (0.021) & (0.008) & (0.010) & (0.017) \\
ln parental income $\times$ Black & $-0.100^{***}$ & $-0.085^{***}$ & $-0.013$ & $-0.011$ & $0.014$ & $0.033^{**}$ \\
 & (0.027) & (0.016) & (0.023) & (0.008) & (0.012) & (0.014) \\
ln parental income $\times$ Black $\times$ age 35--55 & $-0.157^{***}$ & $-0.029^{*}$ & $-0.134^{***}$ & $-0.008$ & $-0.007$ & $-0.007$ \\
 & (0.035) & (0.016) & (0.031) & (0.011) & (0.015) & (0.025) \\
Son-years & 42,238 & 39,704 & 39,704 & 28,754 & 35,251 & 18,660 \\
\addlinespace
\multicolumn{7}{l}{\textit{B. Conditional on education and the wage and occupation score at 28--32}} \\
ln parental income (centred) & $0.006$ & $0.011^{***}$ & $-0.010$ & $0.011^{*}$ & $-0.004$ & $-0.018^{*}$ \\
 & (0.008) & (0.004) & (0.008) & (0.007) & (0.008) & (0.010) \\
\quad $\times$ age 35--55 & $0.138^{***}$ & $-0.002$ & $0.145^{***}$ & $-0.001$ & $0.005$ & $0.030^{*}$ \\
 & (0.021) & (0.010) & (0.021) & (0.009) & (0.010) & (0.016) \\
Black & $-0.015^{*}$ & $-0.026^{***}$ & $0.011$ & $-0.020^{***}$ & $0.015^{*}$ & $-0.009$ \\
 & (0.008) & (0.004) & (0.009) & (0.006) & (0.008) & (0.010) \\
Black $\times$ age 35--55 & $-0.049^{**}$ & $-0.035^{***}$ & $-0.007$ & $-0.003$ & $-0.007$ & $0.040^{**}$ \\
 & (0.020) & (0.011) & (0.020) & (0.008) & (0.011) & (0.019) \\
ln parental income $\times$ Black & $0.023^{**}$ & $-0.006$ & $0.032^{***}$ & $-0.003$ & $0.007$ & $0.039^{***}$ \\
 & (0.010) & (0.005) & (0.011) & (0.008) & (0.012) & (0.014) \\
ln parental income $\times$ Black $\times$ age 35--55 & $-0.117^{***}$ & $-0.005$ & $-0.120^{***}$ & $-0.011$ & $-0.012$ & $-0.002$ \\
 & (0.030) & (0.014) & (0.030) & (0.011) & (0.016) & (0.026) \\
Son-years & 38,565 & 36,904 & 36,904 & 27,010 & 33,203 & 17,231 \\
\bottomrule
\end{tabular}}}
\par\smallskip
\begin{minipage}{\textwidth}\footnotesize \textit{Notes:} Black and White sons pooled (unweighted), son-years at 25--55; one regression per column with the full set of interactions of centred log parental income, a Black indicator and an indicator for ages 35--55; controls: quartic in age, year fixed effects. The triple interaction is the difference between Black and White sons in the after-35 change of the parental-income gradient. Standard errors clustered by 1968 family. Occupational wage score: mean log hourly wage of men aged 25--64 in the occupation relative to all men (CPS MORG, base 1992--2002), occupations harmonised to Dorn's occ1990dd scheme; log wage = score (between occupations) + residual (within occupation), so the score and residual columns add up to the log-wage column on the same sample. $^{***}$ $p<0.01$, $^{**}$ $p<0.05$, $^{*}$ $p<0.1$. Source: \texttt{27.CareerWageGrowth.py} (\texttt{output/career\_wage\_growth.csv}); table built by \texttt{28.PaperBTables.py}.\end{minipage}
\end{table}

\begin{table}[!htbp]
\centering
\caption{Job loss, wages and occupational progression (PSID, sons)}
\label{tab:careerdisp}
\resizebox{\ifdim\width>\textwidth\textwidth\else\width\fi}{!}{%
\resizebox*{!}{\ifdim\totalheight>0.85\textheight 0.85\textheight\else\totalheight\fi}{%
\begin{tabular}{lccccc}
\toprule
Row & \begin{tabular}[b]{@{}c@{}}Change in\\ln wage\end{tabular} & \begin{tabular}[b]{@{}c@{}}Change in\\occupation\\score\end{tabular} & \begin{tabular}[b]{@{}c@{}}Change in\\within-occupation\\residual\end{tabular} & \begin{tabular}[b]{@{}c@{}}Moved up\\$\geq 0.10$\end{tabular} & \begin{tabular}[b]{@{}c@{}}Moved down\\$\geq 0.10$\end{tabular} \\
 & (1) & (2) & (3) & (4) & (5) \\
\midrule
Displaced in $t$: Pooled (weighted) & $-0.131^{***}$ & $-0.028^{***}$ & $-0.103^{***}$ & $0.050^{***}$ & $0.140^{***}$ \\
 & (0.020) & (0.010) & (0.022) & (0.014) & (0.016) \\
Displaced in $t$: Black & $-0.110^{***}$ & $-0.008$ & $-0.103^{***}$ & $0.088^{***}$ & $0.103^{***}$ \\
 & (0.027) & (0.011) & (0.029) & (0.019) & (0.019) \\
Displaced in $t$: White & $-0.148^{***}$ & $-0.035^{***}$ & $-0.112^{***}$ & $0.039^{***}$ & $0.147^{***}$ \\
 & (0.021) & (0.010) & (0.023) & (0.014) & (0.017) \\
Displaced in $t$: White $-$ Black & $-0.037$ & $-0.028^{*}$ & $-0.010$ & $-0.049^{**}$ & $0.045^{*}$ \\
 & (0.034) & (0.015) & (0.037) & (0.024) & (0.025) \\
Son-years (pooled / Black / White) & 23,318 / 6,791 / 16,532 &  &  &  &  \\
\bottomrule
\end{tabular}}}
\par\smallskip
\begin{minipage}{\textwidth}\footnotesize \textit{Notes:} Son-years $t$ (male heads 25--52, 1970--2019) with a valid wage and occupation before ($t-2$ to $t-1$) and after ($t+2$ to $t+3$). Displaced: plant closure, layoff or firing dated $t$; non-displaced son-years with a displacement in $t-3$ to $t+3$ dropped. Outcome = after minus before. Controls: log parental income, quartic in age, year fixed effects. Standard errors clustered by 1968 family. Occupational wage score: mean log hourly wage of men aged 25--64 in the occupation relative to all men (CPS MORG, base 1992--2002), occupations harmonised to Dorn's occ1990dd scheme; log wage = score (between occupations) + residual (within occupation), so the score and residual columns add up to the log-wage column on the same sample. $^{***}$ $p<0.01$, $^{**}$ $p<0.05$, $^{*}$ $p<0.1$. Source: \texttt{IGE/27.CareerWageGrowth.py} (\texttt{output/career\_wage\_growth.csv}); table built by \texttt{IGE/28.PaperBTables.py}.\end{minipage}
\end{table}

\begin{table}[!htbp]
\centering
\caption{Education and instability as mediators of the IGE (PSID, sons)}
\label{tab:mediation}
\resizebox{\ifdim\width>\textwidth\textwidth\else\width\fi}{!}{%
\resizebox*{!}{\ifdim\totalheight>0.85\textheight 0.85\textheight\else\totalheight\fi}{%
\begin{tabular}{lcccccccc}
\toprule
Panel / row & \begin{tabular}[b]{@{}c@{}}ln earnings\end{tabular} & \begin{tabular}[b]{@{}c@{}}ln earnings\end{tabular} & \begin{tabular}[b]{@{}c@{}}ln earnings\end{tabular} & \begin{tabular}[b]{@{}c@{}}ln wage\end{tabular} & \begin{tabular}[b]{@{}c@{}}ln wage\end{tabular} & \begin{tabular}[b]{@{}c@{}}ln hours\end{tabular} & \begin{tabular}[b]{@{}c@{}}ln hours\end{tabular} & \begin{tabular}[b]{@{}c@{}}ln hours\end{tabular} \\
 & (1) & (2) & (3) & (4) & (5) & (6) & (7) & (8) \\
\midrule
\multicolumn{9}{l}{\textit{Black sons}} \\
ln parental family income & $0.214^{***}$ & $0.141^{***}$ & $0.117^{***}$ & $0.164^{***}$ & $0.115^{***}$ & $0.050^{***}$ & $0.027$ & $0.013$ \\
 & (0.039) & (0.036) & (0.029) & (0.029) & (0.026) & (0.018) & (0.019) & (0.015) \\
Son's years of education & No & Yes & Yes & No & Yes & No & Yes & Yes \\
Son's shares of unstable years & No & No & Yes & No & No & No & No & Yes \\
Share of the IGE mediated &  & 0.34 & 0.45 &  & 0.30 &  & 0.46 & 0.75 \\
Child-years & 14,667 & 14,667 & 14,667 & 14,667 & 14,667 & 14,667 & 14,667 & 14,667 \\
\addlinespace
\multicolumn{9}{l}{\textit{White sons}} \\
ln parental family income & $0.447^{***}$ & $0.337^{***}$ & $0.303^{***}$ & $0.413^{***}$ & $0.308^{***}$ & $0.034^{***}$ & $0.029^{**}$ & $0.012$ \\
 & (0.036) & (0.035) & (0.032) & (0.033) & (0.032) & (0.013) & (0.014) & (0.012) \\
Son's years of education & No & Yes & Yes & No & Yes & No & Yes & Yes \\
Son's shares of unstable years & No & No & Yes & No & No & No & No & Yes \\
Share of the IGE mediated &  & 0.25 & 0.32 &  & 0.26 &  & 0.14 & 0.64 \\
Child-years & 27,562 & 27,562 & 27,562 & 27,562 & 27,562 & 27,562 & 27,562 & 27,562 \\
\bottomrule
\end{tabular}}}
\par\smallskip
\begin{minipage}{\textwidth}\footnotesize \textit{Notes:} One regression per column; Lee--Solon design (IGE at 40), sons' child-years at 25--55 with a valid wage. Share mediated $= 1 - \beta_{\text{with controls}}/\beta_{\text{without}}$ relative to the first column of the same outcome. Unstable years: shares of the son's years with unemployment, displacement, reduced hours and zero hours. Education and instability are themselves outcomes of family background, so the shares are an accounting of coefficient movements, not causal mediation. Unweighted; standard errors clustered by 1968 family. $^{***}$ $p<0.01$, $^{**}$ $p<0.05$, $^{*}$ $p<0.1$. Source: \texttt{24.IGEDecomposition.py} (\texttt{output/ige\_decomposition.csv}); table built by \texttt{28.PaperBTables.py}.\end{minipage}
\end{table}

\begin{table}[!htbp]
\centering
\caption{Wage persistence by age, between and within occupations (NLSY79, men)}
\label{tab:nlsycareer}
\resizebox{\ifdim\width>\textwidth\textwidth\else\width\fi}{!}{%
\resizebox*{!}{\ifdim\totalheight>0.85\textheight 0.85\textheight\else\totalheight\fi}{%
\begin{tabular}{lccc}
\toprule
Panel / row & \begin{tabular}[b]{@{}c@{}}ln wage\end{tabular} & \begin{tabular}[b]{@{}c@{}}Occupation\\score\end{tabular} & \begin{tabular}[b]{@{}c@{}}Within-\\occupation\\residual\end{tabular} \\
 & (1) & (2) & (3) \\
\midrule
\multicolumn{4}{l}{\textit{Pooled (weighted)}} \\
ln parental income, ages 25--34 & $0.252^{***}$ & $0.133^{***}$ & $0.119^{***}$ \\
 & (0.014) & (0.008) & (0.012) \\
ln parental income, ages 35--44 & $0.345^{***}$ & $0.128^{***}$ & $0.217^{***}$ \\
 & (0.020) & (0.009) & (0.016) \\
ln parental income, ages 45--55 & $0.381^{***}$ & $0.105^{***}$ & $0.276^{***}$ \\
 & (0.026) & (0.011) & (0.023) \\
\addlinespace
\multicolumn{4}{l}{\textit{Black}} \\
ln parental income, ages 25--34 & $0.161^{***}$ & $0.076^{***}$ & $0.085^{***}$ \\
 & (0.017) & (0.011) & (0.013) \\
ln parental income, ages 35--44 & $0.192^{***}$ & $0.069^{***}$ & $0.123^{***}$ \\
 & (0.023) & (0.013) & (0.019) \\
ln parental income, ages 45--55 & $0.215^{***}$ & $0.061^{***}$ & $0.153^{***}$ \\
 & (0.029) & (0.014) & (0.024) \\
\addlinespace
\multicolumn{4}{l}{\textit{White}} \\
ln parental income, ages 25--34 & $0.241^{***}$ & $0.120^{***}$ & $0.121^{***}$ \\
 & (0.019) & (0.010) & (0.015) \\
ln parental income, ages 35--44 & $0.329^{***}$ & $0.107^{***}$ & $0.222^{***}$ \\
 & (0.025) & (0.011) & (0.020) \\
ln parental income, ages 45--55 & $0.354^{***}$ & $0.086^{***}$ & $0.268^{***}$ \\
 & (0.032) & (0.013) & (0.028) \\
\addlinespace
\multicolumn{4}{l}{\textit{White $-$ Black}} \\
ln parental income, ages 25--34 & $0.080^{***}$ & $0.044^{***}$ & $0.036^{*}$ \\
 & (0.026) & (0.015) & (0.020) \\
ln parental income, ages 35--44 & $0.137^{***}$ & $0.038^{**}$ & $0.099^{***}$ \\
 & (0.034) & (0.017) & (0.027) \\
ln parental income, ages 45--55 & $0.140^{***}$ & $0.025$ & $0.115^{***}$ \\
 & (0.043) & (0.020) & (0.037) \\
Change of the gap, (45--55) $-$ (25--34) & $0.059^{*}$ & $-0.019$ & $0.079^{**}$ \\
 & (0.036) & (0.019) & (0.037) \\
\bottomrule
\end{tabular}}}
\par\smallskip
\begin{minipage}{\textwidth}\footnotesize \textit{Notes:} Men of the NLSY79 living with a parent in 1979; regressor: log mean parental family income 1979--81; outcome: hourly wage of the main (CPS) job at each interview and its occupation score; controls: quartic in age and birth-year fixed effects; standard errors clustered by 1979 household; pooled weighted by the 1979 sampling weight. Occupational wage score: mean log hourly wage of men aged 25--64 in the occupation relative to all men (CPS MORG, base 1992--2002), occupations harmonised to Dorn's occ1990dd scheme; log wage = score (between occupations) + residual (within occupation), so the score and residual columns add up to the log-wage column on the same sample. $^{***}$ $p<0.01$, $^{**}$ $p<0.05$, $^{*}$ $p<0.1$. Source: \texttt{27.CareerWageGrowth.py} (\texttt{output/career\_wage\_growth.csv}); table built by \texttt{28.PaperBTables.py}.\end{minipage}
\end{table}

